\section{Policy Gradient Methods}

\subsection{Stochastic Gradient Decent methods}
\acrfull{sgd} methods are a class of function approximation learning methods well suited for the field of \acrfull{rl}. In these methods the weight vector is a real valued column vector with a fixed number of elements, $\nd{w}=\left (w_1, w_2, ..., w_n \right )^T$, and the approximate function $\hat{v}\left(s, \nd{w}\right)$ is differentiable w.r.t. all weights in $\nd{w}$ at all states $s \in \mathcal{S}$.

The goal is to update the weights in $\nd{w}$ after each observation such that the approximate function $\hat{v}\left(s, \nd{w}\right)$ better approximates the real function $v_{\pi}\left(s\right)$. This can be done by updating the weights in the direction that would reduce the error the most. A definition of this error is thus needed, one possibility is to use the \emph{Mean Squared Value Error} $\overline{\text{VE}}$ shown below \cite{sutton2018reinforcement}.

\begin{equation}
    \overline{\text{VE}}\left(\nd{w}\right) \doteq  \sum_{s\in S} \mu(s) \left[v_{\pi}\left(s\right) - \hat{v}\left(s, \nd{w}\right) \right]^2
\end{equation}

Here is $\mu(s)$ the \emph{state weight distribution}, it is used to indicate which states are more important. This may be necessary since the number of weights in $\nd{w}$ can be significantly smaller than the total number of states, thus prioritizing certain states may be needed.

The following update rule for the weights in $\nd{w}$ may be defined, using $\overline{\text{VE}}$ and assuming the probability for each state to appear is uniformly distributed.

\begin{align}
  \nd{w}_{t+1} &\doteq \nd{w}_t - \frac{1}{2} \alpha \nabla \left[v_{\pi}\left(S_t\right) - \hat{v}\left(S_t, \nd{w}\right) \right]^2 \\
  &= \nd{w}_t + \alpha \left[v_{\pi}\left(S_t\right) - \hat{v}\left(S_t, \nd{w}\right) \right] \nabla \hat{v}\left(S_t, \nd{w}\right)
  \label{eq:w_update}
\end{align}

$\alpha$ is a positive step-size parameter. $\nabla f\left(\nd{w}\right)$ is a vector with the partial derivatives w.r.t. each of the weights. 

\begin{equation}
  \nabla f\left(\nd{w}\right) \doteq \left(\frac{\delta f(\nd{w})}{\delta w_1}, \frac{\delta f(\nd{w})}{\delta w_2}, ..., \frac{\delta f(\nd{w})}{\delta w_n}\right)
\end{equation}

\acrshort{sgd} methods are \emph{Gradient Decent} because the updates to the weight vector $\nd{w}$ are proportional to the negative gradient of each samples squared error. This is the direction where error falls the fastest. Gradient decent methods are called \emph{stochastic} when the updates are done using single samples which may have been selected stochastically. In these cases the updates apply small corrections to the weights and over the course of may iterations it has the effect of reducing the error.

In practice the real value function $v_{\pi}\left(s\right)$ is unknown and thus it's not possible to use Equation~\ref{eq:w_update}. However using Equation~\ref{eq:vpi} it's possible to substitute the value function with the expected return $G_t$. Thus the weight vector update rule can be rewritten as,

\begin{equation}
    \nd{w}_{t+1} = \nd{w}_t + \alpha \left[G_t - \hat{v}\left(S_t, \nd{w}\right) \right] \nabla \hat{v}\left(S_t, \nd{w}\right)
    \label{eq:w_update_2}
  \end{equation}
 
\subsection{Actor-Critic}

\subsection{PPO}
